\section{Radial Distortion Removal}\label{sec:radialdistortion}

In the course, we saw that radial distortion can be modeled as follows:
\begin{equation}\label{eq:radialdistortion}
\begin{cases}
    x'&=p_x+(x-p_x)(1+k_1 r^2+k_2 r^4)\\
    y'&=p_y+(y-p_y)(1+k_1 r^2+k_2 r^4)
\end{cases}
\end{equation}
Where the point $(x,y)$ lands in $(x',y')$ in the distorted image, with
$r^2=\left\| (x,y)-(p_x,p_y) \right\|_2^2=(x-p_x)^2+(y-p_y)^2$.
This means that all points $(x,y)$ at an equal distance $d$ from the principal point $p=(p_x,p_y)$ get distorted the
same amount.
Furthermore, only knowing $p$ and $k_1,k_2$ enables one to invert any distortion by inverting
equation~\eqref{eq:radialdistortion}.

\question{How does the algorithm know that it has converged to a good (or good enough) solution?
What is the error function, exactly?}

\question{Try polynomial and division models. Comment on the difference in results.}

\question{Explain the difference in the principal point position between [A] and [B].}

\question{Explain why we see blurry regions in the output result.}

\question{Is the result reliable? Please comment.}

\question{Show which kind of radial distortion you get with positive and negative k-values.}

Let us assume that $p=(0,0)$ is the center of the image plane to simplify computations and consider the line $y=1$.
We are interested in what happens to the $y$ coordinate of a point $(x,1)$ on the line.
By equation~\eqref{eq:radialdistortion}, we have:

\[y'=1+k_1 (x^2+1)+k_2 (x^2+1)^2\]

Taking the derivative, we obtain
\[\frac{dy'}{dx}=2k_1 x+4k_2 x(x^2+1)\]

If $k_1,k_2<0$, we find that $y'$ is increasing for $x<0$ and decreasing for $x>0$.
Hence, we get barrel distortion.
If $k_1,k_2>0$, we find that $y'$ is decreasing for $x<0$ and increasing for $x<0$.
This is pincushion distortion.

Note that with this model, if $k_1$ is of a different sign than $k_2$, we can have that horizontal lines get distorted
in ways that are neither \enquote*{barrel} nor \enquote*{pincushion} (one can picture some arbitrary fourth-degree
polynomial).